%=============================================================================!
% 16.0  CHAPTER OPENING                                                       !
%=============================================================================!
A trajectory records positions and velocities. To learn from it, we
need to turn those records into quantities that answer physical
questions: how the atoms are arranged, how far they move, or how often
they react. Such an observable may be a single number or a whole curve.
It also gives us something to compare between simulations, and, where
the measurement corresponds, with experiment.

We begin with structure: the distribution of neighbours, the structure
factor measured by diffraction, and the bonds that identify molecules
and reactions. Displacements and velocity correlations then give two
routes to diffusion, followed by the transport of charge and momentum.
The Fourier transform separates vibrations by frequency. Finally,
histograms of selected coordinates give free-energy profiles.

The error estimates of Chapter~15 accompany these measurements. Identities
between observables and comparisons with independent implementations
provide checks on the analysis. The same quantities will later be used
to judge learned models of motion against their reference trajectories
in the planned Part~VI.
